\begin{note}
  \leavevmode
  \begin{equation}
    \label{eq:decomposition}
    (V,T) \cong \qo{\KK[x]}{(f_1)} \oplus \cdots \oplus \qo{\KK[x]}{(f_r)}
  \end{equation}
\end{note}

\begin{remark}
  One can compute the Smith Normal Form effectively using rows and column operations, to use the
  \textbf{division algorithm}, which really require that we're in an Euclidean domain. (\( \KK \)
  field \( \implies \KK[x] \) Euclidean Domain). In particular for PID case there is also a
  generalized version of Smith Normal Form.
\end{remark}

\begin{conclusion}
  \leavevmode
  \begin{enumerate}
    \item Using \textbf{Formula} \ref{eq:decomposition}, we get a basis for \( V \), such that
      the matrix of \( T \) is given as following block decomposition:
      \[
        \begin{pmatrix}
          \cC_{f_1} &           &        & 0         \\
                    & \cC_{f_2} &        &           \\
                    &           & \ddots &           \\
          0         &           &        & \cC_{f_r}
        \end{pmatrix}
      \]
      where \( f_i = x_i^{d_i} + a_{i,1}x^{d_i - 1} + \cdots + a_{i,d_i} \) and we choose the
      basis of \( \qo{\KK[x]}{(f_i)} \) given by \( \overline{1}, \overline{x},\ldots,
      \overline{x}^{d_i - 1} \) and see that
      \[
        \cC_{f_i} =
        \begin{pmatrix}
          0      & 0      &       \cdots & 0      & -a_{i,d_i}     \\
          1      & 0      &       \cdots  & 0      & -a_{i,d_i - 1} \\
          0      & 1      &  \cdots      & \vdots & \vdots         \\
          \vdots & \vdots & \cdots & \vdots & \vdots         \\
          0      & 0      &    \cdots    & 1      & -a_{i,1}
        \end{pmatrix} \text{ for } 1 \leq i \leq r
      \]
    \item See that
      \[
        B_1 \cdot (xI - A ) \cdot B_2 =
        \begin{pmatrix}
          1 &        &   &     &     &        & 0   \\
            & \ddots &   &     &     &        &     \\
            &        & 1 &     &     &        &     \\
            &        &   & f_1 &     &        &     \\
            &        &   &     & f_2 &        &     \\
            &        &   &     &     & \ddots &     \\
          0 &        &   &     &     &        & f_r \\
        \end{pmatrix}
      \]
      take determinant on both sides, we get
      \[
        \underbrace{\det(B_1)}_{\in \KK^{\times}} \cdot \det(xI - A) \cdot \underbrace{\det(B_2)}_{\in \KK^\times}
        = \prod_{i=1}^r f_i
      \]
      which is by the fact that
      \[
        \det (I_n - A) =
        \begin{vmatrix}
          x-a_{11} & - a_{12} & \cdots & -a_{1n}  \\
          -a_{21}  & x-a_{22} & \cdots & -a_{2n}  \\
                   &          & \ddots &          \\
          -a_{n 1} & -a_{n 2} & \cdots & x-a_{nn}
        \end{vmatrix} = x^n + \text{l.o.t.} \implies \text{monic}
      \]
      note that \( f_1,\ldots, f_r \) being monic \( \implies \det(xI - A) = \prod_{i=1}^r f_i \implies
      \det(B_1) \cdot \det(B_2) = 1\).

      \begin{remark}
        Simple fact:
        \[
          c_T (x) = x^n - x^{n-1}\underbrace{(a_{11} + \cdots + a_{nn})}_{\tr(A)} + \cdots +
          (-1)^n \det(A)
        \]
      \end{remark}
  \end{enumerate}
\end{conclusion}

\begin{definition}{Characteristic Polynomial}
  The \underline{\textbf{characteristic polynomial}} of \( T \) is defined as
  \[
    c_T(x) = \det(xI_n - A)
  \]
\end{definition}

\begin{note}
  This is independent of the choice of bases, the matrix for a different basis is \( BAB^{-1} \)
  with \( B \) being invertible, thus
  \[
    \det(xI - BAB^{-1}) = \det(B(xI - A)B^{-1}) = \det(B) \cdot \det(xI - A) \cdot \det(B^{-1}) =
    \det(xI - A)
  \]
\end{note}

\begin{conclusion}
  \leavevmode
  \[
    c_T(x) = \prod_{i=1}^r f_i
  \]
\end{conclusion}
Now consider the set
\[
  \{g \in \KK[x] \; | \; g(T) = 0 \text{ on } V\} = \ann_R(M)
\]
is an ideal in \( \KK[x] \) and recall that as we have \( g(T)(v) = g(x) \cdot v \).

\begin{claim}
  This is \underline{non-zero}, its monic genertator is the \underline{minimal polynomial} \( m_T(x) \)
  of \( T \). \( \iff m_T(x)\)  is the monic polynomial of smallest degree, s.t. \( m_T(T) = 0 \).
\end{claim}
Since
\[
  M \cong \qo{\KK[x]}{(f_1)} \oplus \cdots \oplus \qo{\KK[x]}{(f_r)}
\]
let \( f\in \KK[x] \implies f\cdot (0,\ldots, 0,\overline{1}, 0,\ldots, 0) = 0 \) iff \( f\in (f_i) \; \forall \; i \)
where \( \overline{1} \) is at the i-th position. Hence \( \ann_R(M) = (f_r) \) since \( f_1\mid
f_2\mid \cdots \mid f_r \implies m_T(x) = f_r \).

\begin{conclusion}
  \leavevmode
  \begin{enumerate}
    \item \( m_T \mid c_T \implies c_T(T) = 0 \), which is the \textcolor{blue}{Cayley-Hamilton
      Theorem}.
    \item Since \( f_1\mid f_2 \mid \cdots \mid f_r \implies m_T \) and \( c_T \) has the same
      irreducible factors.
  \end{enumerate}
\end{conclusion}

\begin{note}
  \leavevmode
  \begin{enumerate}
    \item Two similar linear transformation have the same minimal polynomial (the corresponding
      module are isomorphic \( \implies \) they have the same annihilator).
    \item Can define \( m_A, c_T \) when \( A \) is a matrix by taking the correspoinding linear
      transformation.
    \item \( \lambda \in \KK \) is an \underline{eigenvalue} of \( T \) if \( \exists \; v \in V
      \smallsetminus \{0\}\) (called \underline{eigenvector}), s.t. \( T(v) = \lambda v \iff
      (\lambda \Id - T) \) not injective \( \iff \det(\lambda \Id - T) = 0 \iff c_T(\lambda) = 0\).
  \end{enumerate}
\end{note}

\begin{note}[\textbf{Regarding Computation}]
  \leavevmode
  \begin{enumerate}
    \item \( c_T(x) \) is easiest to compute.
    \item If can factor \( c_T \) as product of irreducible polynomials, then can find the minimal
      polynomial by ``trial and error'' (evaluate \( T \)) and get \( 0 \).
    \item Sometimes this is enough (e.g. if \( \deg\bace{\frac{c_T(x)}{m_T(x)}} \leq 1 \)), it's
      easy to calculate.
    \item In general, need to compute the SNF of \( xI -A \) to determine all \( f_i \).
  \end{enumerate}
\end{note}
We now switch our focus to the case when \( \KK \) is an \textbf{\textcolor{blue}{algebraically
closed field}}. Then we obtain the famous \underline{Jordan Canonical Form} of \( T \). Suppose
now \( \KK \) is algebraically closed, e.g. \( \KK = \CC \), then each \( f_i \) can be factored
as a product of \( (x-\lambda_i)^{d_i} \), then
\[
  M \cong \text{direct sum of terms of the form } \qo{\KK[x]}{(x-\lambda)^m}
\]
Let's now consider for such a factor, the basis given by
\[
  v_1 \coloneqq \overline{(x-\lambda)^{m-1}}, v_2 \coloneqq  \overline{(x-\lambda)^{m-2}}, \ldots, v_{m-1}\coloneqq  \overline{(x-\lambda)},
  v_m \coloneqq  \overline{1}
\]
thus the matrix of \( T \) w.r.t. this basis is given by:
\[
  \begin{pmatrix}
    \lambda & 1       &        &        & 0         \\
            & \lambda & \ddots &        &           \\
            &         & \ddots & \ddots &         & \\
            &         &        & \ddots & 1         \\
    0       &         &        &        & \lambda
  \end{pmatrix}
\]
This is a \underline{size-\( m \)} Jordan block correspoinding to \( \lambda \). Which is by the
fact that
\[
  \begin{aligned}
    T(v)   & = x \cdot v                                                                              \\
    T(v_1) & = x \cdot \overline{(x-\lambda^{m-1})}                                                   \\
           & = (x-\lambda + \lambda) \cdot \overline{(x-\lambda)^{m-1}}                               \\
           & = \underbrace{\overline{x(x-\lambda)^{m-1}}}_{=0} + \lambda \overline{(x-\lambda)^{m-1}}
           \\
           & = \lambda\overline{(x-\lambda)^{m-1}} = \lambda v_1                                      \\
    T(v_i) & = x\cdot \overline{(x-\lambda)^{m-i}}                                                    \\
           & = (x-\lambda + \lambda) \cdot \overline{(x-\lambda)^{m-i}}                               \\
           & = \overline{(x-\lambda)^{m-(i-1)}} + \lambda\overline{(x-\lambda)^{m-i}} = v_{i-1} +
           \lambda {v_i}
  \end{aligned}
\]
The \textbf{Upshot}: we get a basis of \( V \), s.t. \( T \) is given by
\[
  \begin{pmatrix}
    J_1 &        &     \\
        & \ddots &     \\
        &        & J_q
  \end{pmatrix}
\]
where \( J_i \) is the Jordan blocks. This is the Jordan canonical form of \( T \) (or \( A \)). The
uniqueness in the structural theorem for modules over PIDs implies that \( \forall \; \lambda,
\forall \; m, \) \# size-\( m \) Jordan blocks correspond to \( \lambda \) is uniquelly determined.
thus we have the uniqueness of the Jordan canonical form up to reordering the blocks.

